\documentclass[12pt,a4paper]{article}

%--------------------------------
% PAQUETES
%--------------------------------
\usepackage{pgfplots}
\usepackage[spanish]{babel}
\usepackage[utf8]{inputenc}
\usepackage[T1]{fontenc}
\usepackage{geometry}
\usepackage{graphicx}
\usepackage{float}
\usepackage{amsmath}
\usepackage{booktabs}
\usepackage{multirow}
\usepackage{enumitem}
\usepackage{fancyhdr}
\usepackage{titlesec}
\usepackage[colorlinks=true, linkcolor=blue, urlcolor=blue]{hyperref}
\usepackage{tcolorbox}
\usepackage{array}
\usepackage{colortbl}
\usepackage{xcolor}

\pgfplotsset{compat=1.18}

\geometry{margin=2.5cm}

\newcommand{\vsample}{\addlinespace}

\setlength{\headheight}{15pt}

\pagestyle{fancy}
\fancyhf{}
\fancyhead[L]{\small Universidad Tecnológica Nacional}
\fancyhead[C]{\small Dispositivos Electrónicos I}
\fancyhead[R]{\small TP N$^\circ$ 3}
\fancyfoot[C]{\thepage}
\renewcommand{\headrulewidth}{0.4pt}
\renewcommand{\footrulewidth}{0pt}

\fancypagestyle{plain}{%
    \fancyhf{}
    \fancyhead[R]{UTN-FRC}
    \fancyhead[C]{Dispositivos Electrónicos I}
    \fancyhead[L]{TP N$^\circ$ 3}
    \fancyfoot[C]{\thepage}
    \renewcommand{\headrulewidth}{0.4pt}
    \renewcommand{\footrulewidth}{0pt}
}

\titleformat{\chapter}[hang]
    {\normalfont\LARGE\bfseries}
    {\thechapter}
    {0.8em}
    {}
\titlespacing*{\chapter}{0pt}{0pt}{1em}

\begin{document}

%--------------------------------
% PORTADA
%--------------------------------
\hypersetup{pageanchor=false}
\begin{titlepage}
    \centering

    {\scshape\LARGE Universidad Tecnológica Nacional \par}
    \vspace{0.3cm}
    {\scshape\Large Facultad Regional Córdoba \par}
    \vspace{1.5cm}

    {\scshape\Large Dispositivos Electrónicos I \par}
    \vspace{2cm}

    {\Huge\bfseries Trabajo Práctico N°3 \par}
    \vspace{0.5cm}
    {\LARGE Transistores Bipolares BJT \par}

    \vspace{2cm}

    \begin{flushleft}
    \textbf{Integrantes:} \\
    Alejos Antony --- 413820 --- Operador\\
    Cortesini Luciano --- 402719 --- Coordinador\\
    Pedregoza Juan --- 408644 --- Documentador\\

    \vspace{0.5cm}

    \textbf{Curso:} 3R3 \\
    \textbf{Grupo:} 02 \\
    \textbf{Docente:}  Leonardo Gabriel Gamboa\\
    \textbf{Año:} 2026 \\
    \end{flushleft}

    \vfill

\end{titlepage}
\hypersetup{pageanchor=true}
\thispagestyle{empty}

\pagestyle{fancy}

\vspace{1cm}

\tableofcontents
\newpage

\section{Introducción}

El presente informe documenta el estudio y la caracterización experimental y simulación del Transistor Bipolar de Juntura (BJT). A lo largo de este
trabajo práctico, se analiza el comportamiento del transistor en sus diferentes regiones de operación (corte, activa y saturación), correlacionando
los modelos teóricos con simulaciones en software (LTspice) y mediciones físicas realizadas en el laboratorio.

El desarrollo del informe se estructura en cinco actividades principales:
\begin{itemize}
    \item \textbf{Actividad 1 (Zona de corte):} Se estudia el transistor con la base abierta para caracterizar y simular
          la corriente de fuga residual ($I_{CEO}$) que circula a través del colector cuando el dispositivo se encuentra en zona de corte.
    \item \textbf{Actividad 2 (Polarización de la juntura BE):} Se analiza el comportamiento exponencial de la unión
          base-emisor ($I_B = f(V_{BE})$) y las  influencia de la temperatura, contrastando los resultados de simulación
          con las mediciones de laboratorio en circuitos de polarización en emisor común.
    \item \textbf{Actividad 3 (Curvas características):} Se obtienen y analizan las curvas de salida del transistor
          ($I_C = f(V_{CE})$) para diferentes niveles de corriente de base ($I_B$), delimitando las fronteras entre las
          zonas de corte, activa y saturación, y estimando la ganancia de corriente ($\beta$).
    \item \textbf{Actividad 4 (Transferencia de corriente):} Se evalúa la relación de transferencia de corriente ($I_C =
          f(I_B)$) a distintos voltajes de colector-emisor ($V_{CE}$) y se estudia la dependencia del parámetro de
          ganancia ($\beta$) en función del punto de operación.
    \item \textbf{Actividad 5 (Especificaciones del fabricante):} Se realiza un análisis comparativo de parámetros
          eléctricos, térmicos y de conmutación de seis modelos comerciales de transistores BJT (tanto NPN como PNP) a
          partir de sus hojas de datos.
\end{itemize}

El objetivo final es lograr una comprensión sólida del funcionamiento práctico del BJT, identificando sus no
linealidades, efectos térmicos y limitaciones de diseño en circuitos electrónicos.
\newpage

\section{Actividad 1: Zona de corte}
En esta actividad se estudia el comportamiento del transistor bipolar (BC547C) en configuración de emisor común cuando
la base se encuentra en circuito abierto, lo que impone teóricamente una corriente de base $I_B = 0\,\mathrm{A}$. El
propósito es verificar el estado de corte del dispositivo y analizar la existencia de la corriente de fuga
colector-emisor con base abierta ($I_{CEO}$). Dado que el orden de magnitud de esta corriente es extremadamente pequeño
(nanoampere) y produce caídas de tensión en el resistor de colector que están por debajo del límite
de resolución de los multímetros digitales comunes, este análisis se lleva a cabo únicamente a través de simulación en
LTspice.

\subsection{Simulación}
Se diseña el circuito esquemático conectando un resistor de colector $R_C = 1\,\mathrm{k}\Omega$ alimentado por una
fuente $V_{CC} = 10\,\mathrm{V}$, y un resistor de base $R_B = 10\,\mathrm{k}\Omega$ con su terminal izquierdo en el
aire. Mediante un análisis de punto de operación en corriente continua (\texttt{.op}), se determina la corriente
residual de colector $I_C = I_{CEO}$ y la tensión colector-emisor $V_{CE}$.

\begin{figure}[H]
    \centering
    \includegraphics[width=1\textwidth]{sim1.png}
    \caption{Simulación del punto de operación en continua (.op) para el transistor en corte (base abierta).}
\end{figure}

Conclusión: 
La simulación DC op point confirma que el transistor BC547C se encuentra en zona de corte: las corrientes de base,
colector y emisor son prácticamente nulas, valores que corresponden únicamente a la corriente de fuga del dispositivo.
Además, la resistencia no presenta caída de tensión y el transistor se comporta como un circuito abierto entre colector
y emisor, confirmando que se encuentra en zona de corte.


\section{Actividad 2: Polarización de la juntura BE}
El propósito de esta actividad es caracterizar el comportamiento eléctrico de la unión base-emisor ($BE$) del transistor
BJT y analizar su sensibilidad frente a variaciones de temperatura. Dado que la unión $BE$ constituye una juntura PN, su
comportamiento es análogo al de un diodo semiconductor. Se realiza primero una caracterización por simulación y luego se
efectúa la correspondiente verificación experimental en el laboratorio.

\subsection{Simulación}
Para modelar la curva de la juntura, se diseña en LTspice un circuito de polarización con $R_C = 560\,\Omega$ y $R_B =
10\,\mathrm{k}\Omega$. Se define una fuente fija $V_{CC} = 10\,\mathrm{V}$ y una fuente de base variable $V_{BB}$.
Mediante un barrido de tensión continua (DC Sweep) de $V_{BB}$ de 0 a 10 V (con pasos de 1 mV), se grafica la corriente
de base $I_B$ en función de la tensión interna base-emisor $V_{BE}$. Posteriormente, se incorpora la directiva
\texttt{.step temp 0 50 25} para evaluar el comportamiento térmico superponiendo los resultados a
$0\,^{\circ}\mathrm{C}$, $25\,^{\circ}\mathrm{C}$ y $50\,^{\circ}\mathrm{C}$.

\begin{figure}[H]
  \centering
  \includegraphics[width=1\textwidth]{Ib(Vbe).png}
  \caption{Curva de transferencia de entrada $I_B = f(V_{BE})$ obtenida mediante simulación.}
\end{figure}

\begin{figure}[H]
  \centering
  \includegraphics[width=1\textwidth]{2Ib(Vbe)-(temperatura).png}
  \caption{Desplazamiento térmico de la característica de entrada $I_B = f(V_{BE})$ para $0\,^{\circ}\mathrm{C}$,
  $25\,^{\circ}\mathrm{C}$ y $50\,^{\circ}\mathrm{C}$.}
\end{figure}
\subsection{Conclusiones}

El barrido de la tensión de base (\(V_{BB}\)) permitió observar el comportamiento característico de la unión
base-emisor. Para tensiones inferiores a aproximadamente \(0.7\,\mathrm{V}\), la corriente de base (\(I_B\)) es
prácticamente nula debido a que la unión no se encuentra polarizada en directa. A partir de dicho umbral, \(I_B\)
aumenta rápidamente con pequeños incrementos de \(V_{BB}\), siguiendo una respuesta exponencial similar a la de un diodo
de silicio, ya que la unión base-emisor presenta el mismo comportamiento eléctrico.

En el barrido de temperatura (\(0\,^{\circ}\mathrm{C}\), \(25\,^{\circ}\mathrm{C}\) y \(50\,^{\circ}\mathrm{C}\)) se
observa un desplazamiento de la curva \(I_B\)-\(V_{BB}\). A medida que aumenta la temperatura, la tensión necesaria para
polarizar la unión base-emisor disminuye, por lo que para un mismo valor de \(V_{BB}\) la corriente de base es mayor.
Esto se debe a la reducción de la barrera de potencial de la unión PN con la temperatura.

En consecuencia, el incremento de la temperatura modifica el punto de operación del transistor, pudiendo provocar
variaciones en las corrientes y tensiones de polarización. Este efecto debe considerarse durante el diseño de circuitos
de polarización para garantizar un funcionamiento estable frente a cambios térmicos.


\subsection{Implementación en el laboratorio}
Se implementa en protoboard el circuito de emisor común con resistores $R_B = 10\,\mathrm{k}\Omega$ y $R_C =
560\,\Omega$ medidos previamente. La fuente $V_{CC}$ se fija en $10\,\mathrm{V}$ y la tensión $V_{BB}$ se
incrementa de forma manual y progresiva desde 0 V. En cada punto establecido se registran las tensiones $V_{BE}$ y la
caída de tensión en el resistor de base $V_{R_B}$. Con esta última se calcula la corriente de base $I_B = V_{R_B} / R_B$
por ley de Ohm. Asimismo, se mide la corriente de colector $I_C$ e indirectamente se determinan las potencias disipadas
en el resistor de colector ($P_{R_C} = I_C^2 \cdot R_C$) y en el transistor ($P_{\mathrm{bjt}} \approx V_{CE} \cdot
I_C$).

\begin{figure}[H]
    \centering
    \includegraphics[width=0.5\textwidth]{ciruictoact2.jpeg}
    \caption{Montaje experimental del circuito de polarización para la juntura base-emisor.}
\end{figure}

\begin{table}[H]
\centering
\caption{Resultados de corrientes y potencias - Actividad 2}
\begin{tabular}{cccccccc}
\toprule
\textbf{$V_{bb}$}  & \textbf{$V_{be}$}  & \textbf{$V_{rb}$}  & \textbf{$I_b$ calculado}  & \textbf{$I_c$} &
    \textbf{$P_{rc}$}  & \textbf{$P_{bjt}$} \\ \midrule
467 mV & 455 mV & tiende a cero & $\approx 0\ \mu\text{A}$ & $\approx 0\text{ mA}$ & 0 mW & 0 mW \\ \vsample
1030 mV & 955 mV & 84 mV & 8,62 $\mu\text{A}$ & 1,29 mA* & 0,92 mW & 12,0 mW \\ \vsample
2,01 V & 0,93 V & 1,07 V & 109,7 $\mu\text{A}$ & 16,45 mA* & 149,1 mW & 15,4 mW \\ \vsample
3,00 V & 0,97 V & 1,95 V & 200,0 $\mu\text{A}$ & 16,8 mA (lab) & 155,5 mW & 12,5 mW \\ \vsample
4,0 V & 0,87 V & 2,90 V & 297,4 $\mu\text{A}$ & 16,8 mA (lab) & 155,5 mW & 12,7 mW \\ \vsample
5,0 V & 0,95 V & 3,95 V & 405,1 $\mu\text{A}$ & 16,8 mA (lab) & 155,5 mW & 12,8 mW \\ \bottomrule
\end{tabular}
\label{tab:actividad2}
\end{table}
\\

\subsection{Análisis}

Para \(V_{BB}=467\,\text{mV}\), la tensión \(V_{BE}\) es insuficiente para polarizar la unión base-emisor en directa,
por lo que el transistor permanece en la región de corte.

Al aumentar \(V_{BB}\) hasta aproximadamente \(1.03\,\text{V}\), el transistor entra en la región activa. En esta zona
la corriente de colector \(I_C\) aumenta proporcionalmente a la corriente de base \(I_B\), de acuerdo con la ganancia de
corriente del transistor.

Cuando \(V_{BB}\) supera aproximadamente los \(2\,\text{V}\), la corriente de colector deja de incrementarse y se
estabiliza alrededor de \(16.8\,\text{mA}\), aun cuando la corriente de base continúa aumentando. Este comportamiento
indica que el transistor ha alcanzado la región de saturación, donde la corriente de colector ya no depende de \(I_B\),
sino que queda limitada por la fuente de alimentación y la resistencia de colector.

La corriente de colector en saturación puede calcularse aplicando la Ley de Kirchhoff a la malla del colector:

\begin{equation}
V_{CC}=I_{C(\mathrm{sat})}R_C+V_{CE(\mathrm{sat})}
\end{equation}

Despejando la corriente de saturación se obtiene:

\begin{equation}
I_{C(\mathrm{sat})}
=\frac{V_{CC}-V_{CE(\mathrm{sat})}}{R_C}
=\frac{10\,\mathrm{V}-0.2\,\mathrm{V}}{560\,\Omega}
=\frac{9.8\,\mathrm{V}}{560\,\Omega}
=0.0175\,\mathrm{A}
\approx17.5\,\mathrm{mA}.
\end{equation}

El valor calculado resulta muy próximo al valor medido experimentalmente (\(I_C \approx 16.8\,\text{mA}\)), lo que
confirma que el transistor se encuentra en saturación y que la corriente de colector está limitada por la resistencia
\(R_C\) y la tensión de alimentación \(V_{CC}\).

\section{Actividad 3: Curvas características}
El objetivo de esta actividad es obtener la familia de curvas características de salida del transistor ($I_C =
f(V_{CE})$) para diferentes valores preestablecidos de corriente de base $I_B$. Estas curvas son fundamentales para
identificar gráficamente las zonas de corte, activa y saturación del transistor BJT, así como para estimar la ganancia
de corriente $\beta$ en la región de amplificación.

\subsection{Simulación}
Se implementa la topología de emisor común con $R_B = 10\,\mathrm{k}\Omega$ y $R_C = 560\,\Omega$. Se define una directiva
de barrido bidimensional (DC Sweep): un barrido principal en la tensión de colector $V_{CC}$ de 0 a 10 V con
incrementos de 0.1 V (para variar $V_{CE}$), y un barrido secundario en la tensión de base $V_{BB}$ de 0 a 5 V con pasos
de 0.2 V (para obtener múltiples trazas asociadas a diferentes corrientes de base $I_B$). Se grafica la corriente de
colector $I_C$ en función de la tensión del nodo de colector $V_C$.

\begin{figure}[H]
    \centering
    \includegraphics[width=1\textwidth]{3Familia_curvas_Ic(Vbb).png}
    \caption{Familia de curvas de salida $I_C = f(V_{CE})$ obtenidas mediante simulación en LTspice.}
\end{figure}

\subsection{Implementación en el laboratorio}
Sobre el circuito montado en protoboard, se realiza un relevamiento manual y cíclico de datos. Para cada escalón de
corriente de base $I_B$ ($0\,\mu\mathrm{A}$, $10\,\mu\mathrm{A}$, $15\,\mu\mathrm{A}$, $20\,\mu\mathrm{A}$ y
$25\,\mu\mathrm{A}$), controlado mediante el ajuste de la fuente $V_{BB}$ de entrada, se realiza un barrido de la fuente
$V_{CC}$ midiendo en cada punto la corriente de colector $I_C$ y el voltaje colector-emisor $V_{CE}$. En caso de no
disponer de suficiente resolución para medir corrientes pequeñas de base, se calcula $I_B$ de forma indirecta midiendo
la caída en $R_B$.

\begin{table}[H]
\centering
\small
\caption{Valores experimentales relevados para la gráfica de curvas de salida ($I_C = f(V_{CE})$).}
\label{tab:curvas_salida_v2}
\begin{tabular}{lcccccccc}
\toprule
\textbf{Parámetro} & \multicolumn{7}{c}{\textbf{Voltaje de Colector $V_{CC}$ [V]}} \\
\cmidrule{2-8}
 & \textbf{0} & \textbf{0,25} & \textbf{0,5} & \textbf{1} & \textbf{2} & \textbf{5} & \textbf{10} \\ \midrule
\textbf{$I_B = 0\ \mu\text{A}$} & & & & & & & \\
$I_C$ [mA] & 0 & 0 & 0 & 0 & 0 & 0 & 0 \\
$V_{CE}$ [V] & 0 & 0,25 & 0,5 & 1 & 2 & 5 & 10 \\ \midrule
\textbf{$I_B = 10\ \mu\text{A}$} & & & & & & & \\
$I_C$ [mA] & 0,02 & 0,20 & 0,41 & 0,97 & 1,42 & 1,48 & 1,50 \\
$V_{CE}$ [V] & 0,042 & 0,046 & 0,095 & 0,162 & 0,630 & 3,720 & 8,800 \\ \midrule
\textbf{$I_B = 15\ \mu\text{A}$} & & & & & & & \\
$I_C$ [mA] & 0,04 & 0,45 & 0,48 & 1,04 & 2,20 & 2,55 & 2,68 \\
$V_{CE}$ [V] & 0,050 & 0,085 & 0,090 & 0,140 & 1,200 & 2,780 & 7,850 \\ \midrule
\textbf{$I_B = 20\ \mu\text{A}$} & & & & & & & \\
$I_C$ [mA] & 0 & 0,50 & 0,55 & 1,00 & 2,10 & 3,15 & 3,75 \\
$V_{CE}$ [V] & 0 & 0,085 & 0,090 & 0,120 & 1,400 & 2,130 & 6,800 \\ \midrule
\textbf{$I_B = 25\ \mu\text{A}$} & & & & & & & \\
$I_C$ [mA] & 0 & 0,51 & 0,55 & 1,21 & 2,10 & 4,53 & 4,70 \\
$V_{CE}$ [V] & 0 & 0,080 & 0,080 & 0,120 & 1,200 & 1,290 & 6,080 \\
\bottomrule
\end{tabular}
\end{table}

\subsection{Análisis}

Para valores mayores de $I_B$, se distinguen dos regiones: una zona de saturación para $V_{CE}$ bajos, donde $I_C$ crece
rápidamente con $V_{CE}$; y una zona activa donde $I_C$ se mantiene aproximadamente constante e independiente de
$V_{CE}$, controlada exclusivamente por $I_B$. A mayor $I_B$, mayor es el valor de $I_C$ en la zona activa, con una
relación aproximada $I_C = \beta \cdot I_B$. El espaciado entre curvas para distintos valores de $I_B$ permite estimar
el $\beta$ del transistor: tomando $I_B = 20\,\mu\text{A}$ e $I_C = 3.75\,\text{mA}$ a $V_{CE} = 10\,\text{V}$, se
obtiene $\beta \approx 187$, consistente con el rango declarado por el fabricante para el BC547C.

\section{Actividad 4: Transferencia de corriente }
En esta actividad se evalúa de manera cuantitativa la relación de transferencia de corriente $I_C = f(I_B)$ y el
comportamiento de la ganancia de corriente continua $\beta$ (o $h_{FE}$) como función de las condiciones de polarización
del circuito. El objetivo es analizar que la ganancia no es un parámetro estático, sino que presenta variaciones ante la
magnitud de la corriente y ante la tensión de salida $V_{CE}$ (efecto de modulación de ancho de base o efecto Early).

\subsection{Simulación}
En la simulación en LTspice, se configura un barrido de continua de la fuente $V_{BB}$ (0 a 5 V con pasos de 0.01 V)
como barrido principal, y un barrido tipo lista para la fuente de colector $V_{CC}$ (con valores de 2, 5 y 8 V) como
barrido secundario. Se grafican sucesivamente:
\begin{enumerate}
    \item La corriente de colector $I_C$ en función de la corriente de base $I_B$ para observar la linealidad de la
          transferencia.
    \item El cociente matemático $I_C / I_B$ en función de $I_B$ para obtener la curva de ganancia instantánea $\beta$.
    \item La ganancia $\beta$ en función de la corriente de colector $I_C$, identificando la meseta de máxima ganancia y
          la caída en los extremos de operación.
\end{enumerate}

\begin{figure}[H]
    \centering
    \includegraphics[width=1\textwidth]{sim41.png}
    \caption{Curvas de transferencia de corriente $I_C = f(I_B)$ para $V_{CE}$ de 2 V, 5 V y 8 V.}
\end{figure}
\begin{figure}[H]
    \centering
    \includegraphics[width=1\textwidth]{sim42.png}
    \caption{Ganancia de corriente instantánea $\beta$ en función de la corriente de base $I_B$.}
\end{figure}
\begin{figure}[H]
    \centering
    \includegraphics[width=1\textwidth]{sim43.png}
    \caption{Variación de la ganancia de corriente $\beta$ en función de la corriente de colector $I_C$ bajo diferentes
    tensiones $V_{CE}$.}
\end{figure}

\subsection{Implementación en el laboratorio}
Se utiliza el circuito en emisor común alimentado con tensión $V_{CC}$ variable. Se varía la fuente $V_{BB}$ para fijar
corrientes de base específicas de $5\,\mu\mathrm{A}$, $7\,\mu\mathrm{A}$, $10\,\mu\mathrm{A}$ y $20\,\mu\mathrm{A}$.
Puesto que la caída de tensión en el resistor de colector modifica la tensión $V_{CE}$ real al variar la corriente, en
cada medición se reajusta levemente el valor de $V_{CC}$ para asegurar que $V_{CE}$ se mantenga constante en los tres
niveles de ensayo: 2 V, 5 V y 8 V.

\begin{table}[H]
\centering
\caption{Valores experimentales relevados para la transferencia de corriente en función de $V_{CE}$.}
\label{tab:transferencia_corriente}
\begin{tabular}{cccc}
\toprule
\textbf{$I_B$ [$\mu$A]}  & \textbf{$I_C$ [mA]} (@$V_{CE} = 2$V)  & \textbf{$I_C$ [mA]} (@$V_{CE} = 5$V) &
    \textbf{$I_C$ [mA]} (@$V_{CE} = 8$V) \\ \midrule
5  & 0,90 & 0,92 & 0,933 \\ \vsample
7  & 1,46 & 1,50 & 1,520 \\ \vsample
10 & 2,26 & 2,29 & 2,330 \\ \vsample
20 & 5,00 & 5,15 & 5,180 \\ \bottomrule
\end{tabular}
\end{table}

\subsection{Análisis}

Aquí se muestran las características de transferencia $I_C = f(I_B)$ para distintos valores de $V_{CE}$. En todos los
casos la relación es lineal en la zona activa, con una pendiente que corresponde a la ganancia de corriente $\beta$. Las
curvas para los tres valores de $V_{CE}$ son prácticamente idénticas, lo que confirma que en la zona activa la corriente
de colector $I_C$ es muy poco sensible a las variaciones de $V_{CE}$.

\section{Actividad 5: Especificaciones del fabricante }
Esta actividad tiene como finalidad familiarizarse con la lectura e interpretación de las hojas de datos (datasheets)
provistas por los fabricantes de componentes semiconductores. A partir de ellas, se identifican y comparan los límites
máximos absolutos, las características térmicas y los tiempos de respuesta de conmutación para diversos modelos de
transistores bipolares.

\subsection{Análisis}
Se descargan las hojas de datos oficiales y se releva la información para seis transistores bipolares comerciales: BC548
(NPN de uso general), BC557 (PNP complementario de uso general), 2N2222 (NPN de conmutación rápida), BU208 (NPN de alta
tensión), MPS6514 (NPN de pequeña señal) y TIP36 (PNP de potencia). Los parámetros analizados comprenden tensiones de
ruptura, corrientes máximas de operación, ganancias de corriente en continua y alterna ($h_{FE}$ y $h_{fe}$), tensiones
de saturación, potencia disipable máxima ($P_d$), resistencias térmicas ($\theta_{jc}$, $\theta_{ja}$) y tiempo de
encendido ($t_{\mathrm{on}}$).

\begin{table}[H]
\centering
\small
\caption{Relevamiento completo de parámetros críticos extraídos de las hojas de datos.}
\label{tab:datasheets_bjt_completa}
\begin{tabular}{lcccccc}
\toprule
\textbf{Parámetro}  & \textbf{BC548}  & \textbf{BC557}  & \textbf{2N2222}  & \textbf{BU208}  & \textbf{MPS6514} &
    \textbf{TIP36} \\
 & (NPN) & (PNP) & (NPN) & (NPN) & (NPN) & (PNP) \\ \midrule
\multicolumn{7}{l}{\textit{Características Eléctricas Máximas y de Operación (Tj = 25 °C)}} \\ \vsample
$V_{(BR)CEO}$ [V] & 30 V & 45 V & 30 V & 700 V & 25 V & 60 V \\ \vsample
$I_{CEO}$ & 15 nA & 15 nA & 10 nA & 1 mA & 0,5 $\mu$A & 1 mA \\ \vsample
$I_{CES}$ & 10 nA & 100 nA & 10 nA & 1 mA & 0,1 $\mu$A & 2 mA \\ \vsample
$I_C$ max & 100 mA & 100 mA & 800 mA & 8 A & 150 mA & 25 A \\ \vsample
$I_{EBO}$ & 10 nA & 10 nA & 10 nA & 10 mA & 0,1 $\mu$A & 1 mA \\ \vsample
$h_{FE}$ & 110 / 800 & 75 / 800 & 100 / 300 & 2,25 (mín) & 150 / 300 & 15 / 50 \\ \vsample
$h_{fe}$ & 125 / 900 & 100 / 900 & 50 / 300 & 2,50 (mín) & - & 20 (mín) \\ \vsample
$V_{BE}$ [V] & 0,55 V & 0,65 V & 0,6 V & 1,5 V & 0,65 V & 1,1 V \\ \vsample
$V_{CE(sat)}$ [V] & 0,25 V & 0,3 V & 0,3 V & 5 V & 0,25 V & 1,8 V \\ \vsample
$P_d$ [W] & 500 mW & 500 mW & 500 mW & 12,5 W & 625 mW & 125 W \\ \midrule
\multicolumn{7}{l}{\textit{Características Térmicas}} \\ \vsample
$\theta_{jc}$ [$^{\circ}$C/W] & 85 & 80 & 83,3 & 10 & 83,3 & 1,0 \\ \vsample
$\theta_{ja}$ [$^{\circ}$C/W] & 200 & 200 & 300 & 70 & 200 & 35,7 \\ \midrule
\multicolumn{7}{l}{\textit{Características de Conmutación}} \\ \vsample
$t_{on}$ & 35 ns & 50 ns & 35 ns & 0,7 $\mu$s & - & 40 ns \\ \bottomrule
\end{tabular}
\end{table}

Análisis de parámetros: \\
El relevamiento de los parámetros de los seis transistores permite identificar patrones claros según el tipo de
aplicación. En cuanto a la tensión de ruptura $V_{(BR)CEO}$, el BU208 destaca ampliamente con $700\,\text{V}$, debido a
sus aplicaciones para alta tensión a diferencia de los demás modelos. La corriente máxima de colector ($I_C$) refleja la
misma división: los dispositivos de potencia (TIP36: $25\,\text{A}$, BU208: $8\,\text{A}$) superan ampliamente a los de
señal pequeña (BC548 y BC557: $100\,\text{mA}$, 2N2222: $800\,\text{mA}$). \\
La ganancia de corriente en continua $h_{FE}$ muestra una relación inversa con la potencia manejada. \\
Desde el punto de vista térmico, la resistencia térmica juntura-ambiente ($\theta_{JA}$) es mucho menor en los
dispositivos de potencia.
\\
En conmutación, los transistores de señal pequeña presentan tiempos de encendido de $35$ a $50\,\text{ns}$, adecuados
para aplicaciones de alta frecuencia, mientras que el BU208 con $0.7\,\mu\text{s}$ queda limitado a aplicaciones de baja
frecuencia como etapas de salida de deflexión o fuentes de alimentación conmutadas.
\section{Conclusiones}

\begin{itemize}

\item Se comprobó que el transistor bipolar no presenta un comportamiento completamente ideal, ya que aun en la región
      de corte existe una pequeña corriente de fuga (\textit{ICEO}) entre colector y emisor.
\item La comparación entre los valores de \textit{ICEO} obtenidos mediante simulación y los especificados en el
      datasheet mostró diferencias esperables, atribuibles a que los modelos de simulación representan un dispositivo
      típico, mientras que las hojas de datos indican valores máximos garantizados bajo condiciones específicas de
      ensayo. Además, influyen las tolerancias de fabricación y las condiciones de temperatura.

\item El relevamiento experimental permitió observar que la corriente de base ($I_B$) presenta un comportamiento
      exponencial respecto del voltaje $V_{BE}$, identificándose un claro ``codo de conducción'' alrededor de
      $0,6$--$0,7\,\mathrm{V}$, a partir del cual el transistor comienza a conducir de manera significativa.

\item Una vez polarizada en directa la juntura base-emisor, el voltaje $V_{BE}$ permaneció aproximadamente constante,
      variando sólo algunos milivoltios frente a incrementos importantes de corriente. Este comportamiento es
      equivalente al de una juntura PN, por lo que resulta análogo al de un diodo de silicio.

\item El análisis de las simulaciones con variaciones de temperatura mostró que un aumento de ésta provoca una
      disminución del voltaje $V_{BE}$ y un incremento de la corriente de base y de la corriente de colector para una
      misma polarización.

\item Tanto en las mediciones experimentales como en las simulaciones fue posible identificar claramente las tres
      regiones de funcionamiento del transistor: corte, región activa y saturación, verificándose la concordancia entre
      el comportamiento observado y el modelo teórico.

\item La curva de transferencia $I_C=f(I_B)$ presentó una relación prácticamente lineal dentro de la región activa,
      verificándose que la corriente de colector aumenta casi proporcionalmente con la corriente de base. Sin embargo,
      al alcanzarse la región de saturación, $I_C$ dejó de incrementarse a pesar de continuar aumentando $I_B$,
      permaneciendo prácticamente constante en un valor cercano a la corriente de saturación ($I_{C(\mathrm{sat})}$),
      determinada por las condiciones del circuito.

\item Se verificó que la ganancia de corriente continua ($\beta$ o $h_{FE}$) no permanece constante, sino que depende
      del punto de operación, particularmente de la corriente de colector, la tensión $V_{CE}$, la temperatura y las
      características propias del dispositivo. Esto confirma que $\beta$ debe considerarse un parámetro variable en el
      diseño práctico.

\item Se observaron diferencias entre los resultados experimentales, y las simulaciones ,que pueden atribuirse a las
      tolerancias de fabricación del transistor, a las incertidumbres propias de las mediciones y a las aproximaciones
      empleadas en los modelos de simulación.

\end{itemize}




\section{Bibliografía / Referencias}

\begin{itemize}
    \item Floyd, Thomas L.\ \textit{Principles of Electric Circuits: Conventional
    Current Version}. Pearson Prentice Hall, 2007.
 
    \item Gonz\'alez, M\'onica Liliana. \textit{LTSPICE An\'alisis de circuitos y
    dispositivos electr\'onicos}, 1\textsuperscript{a} ed.\ La Plata, EDULP, 2018.
    Disponible en: \url{http://sedici.unlp.edu.ar/handle/10915/69818}.
 
\end{itemize}
 

\end{document}
